\section{Continuous design, correlations, and the direction of optimization}
\label{sec:main-design}

With fixed nominations, every cactus cycle has one scalar circulation.
This permits considerably more coefficient dependence than the independent
edge boxes used in nomination optimization. It also exposes a boundary
between correlations within a cycle and correlations across cycles.

For each edge let
$g_e(x;\theta)=g_{e0}(x)+\sum_j\theta_jg_{ej}(x)$,
with densely encoded rational polynomial pieces, fixed rational
breakpoints, and admissible continuous strictly increasing complete
laws vanishing at zero throughout the parameter domain. The degree,
number of pieces, and parameter dimension may grow. Admissibility is
a promise in this correlated model. Bridge flows are fixed by the
nomination; a consistently oriented cycle has
$x_e=x_e^0+\chi_eq$ and loop equation
\[
 H_C(q,\theta)=\sum_{e\in C}\chi_eg_e(x_e^0+\chi_eq;\theta)=0.
\]
This function is continuous, strictly increasing in $q$, and affine
in $\theta$. Rational bounds on total injection supply a root bracket.

\begin{theorem}[Cycle-local design and global capacity recovery]
\label{thm:main-design}
Let $G$ be a connected simple cactus with fixed rational balanced
nominations and the admissible affine polynomial-law model above.
\begin{enumerate}
\item If each cycle has its own independent nonempty bounded rational
parameter polytope, the attainable flow region is
$x=x^0+Zq$ with independent algebraic circulation intervals.
Rational cycle-local signed capacities clip these intervals exactly.
Feasibility, each interval endpoint, and a rational original parameter
profile attaining it are computable in polynomial bit time. Linear
flow objectives admit exactly optimizing rational profiles and additive
scalar values, with no exact-comparison claim for a growing sum.
\item In the same independent-cycle model with those capacities, a
densely encoded rational polynomial flow objective promised convex on
$[-B-1,B+1]^m$, $B=\sum_v|b_v|$, admits an exactly feasible rational
parameter scenario within $\eps$ of its minimum in polynomial
accuracy-bit time. Degrees and number of cycles are unrestricted.
For maximization of $g(Rx)$, the analogous additive scenario guarantee
holds when $R\in\Q^{k\times m}$ has fixed row count $k$ and the dense
rational polynomial $g$ is convex on a box containing the attainable
measurements.
\item If one nonempty bounded rational parameter polytope instead
couples all cycles, feasibility of signed arc capacities is rational
linear feasibility and has a rational witness. Robust capacity testing
uses rational LPs and returns rational violating profiles when needed.
Individual arcs have exact algebraic extrema and rational optimizing
profiles, including when other such capacities filter the domain.
\end{enumerate}
All complexity bounds count dense input and output bits. These statements
exclude variable nominations and potential filters; part 2 excludes
flow constraints coupling distinct cycle coordinates.
\end{theorem}

These are \cref{thm:a-design-independent,thm:a-design-convex,thm:a-design-measurements,thm:a-design-global-capacity}.
The few-measurement maximization conclusion also holds for independent
finite quadratic resistance sets at fixed nominations, without capacity
filters. Its dependence is $N^{O(k)}$ times a polynomial in accuracy
bits, and it promises a near-optimal scenario rather than exact
comparison of candidate values. A positive-semidefinite quadratic of
fixed rank $r$ uses at most $r+1$ measurements, including its linear
term (\cref{cor:a-design-fixed-rank}).

\subsection{Why scalar roots admit exact rational parameter recovery}

At a rational $z$, strict increase gives
\[
 \max_{\theta\in P}q(\theta)\ge z
 \quad\Longleftrightarrow\quad
 \min_{\theta\in P}H_C(z,\theta)\le0.
\]
The right side is an LP, regardless of the parameter dimension.
A largest root is attained at a vertex of $P$: expressing any profile
as a convex combination of vertices expresses its zero loop value as
the same combination, so at least one vertex has a root no smaller.
Vertex coordinates have polynomial bit bounds from determinants.
After substitution, all vertex-piece roots therefore satisfy a common
degree and coefficient-height bound. Root separation gives a polynomial
number of bisection bits sufficient to distinguish unequal candidate
roots, without enumerating vertices. An LP vertex at the final threshold
must attain the exact extreme root. \Cref{thm:a-design-root} proves this
step for general monotone dense piecewise-polynomial root families,
including repeated roots, breakpoint roots, and lower-dimensional
polytopes.

Capacities are equally direct. A lower bound $q\ge a$ is equivalent
to $H_C(a,\theta)\le0$, and an upper bound $q\le b$ is equivalent
to $H_C(b,\theta)\ge0$. Both are linear in the original parameters.
This remains true under global correlations. Independence across cycles
is needed only to combine the scalar attainable intervals into a box.
A rational circulation target can be realized by another rational LP.
Together with exact endpoint profiles, this permits recovery even when
a clipped interval is an irrational singleton.

Dense convex minimization uses rational inner boxes approximating the
algebraic circulation box, a controlled convex extension for a polynomial
promised convex only on the stated cube, and exact rational oracles.
The optimizer is mapped back to original coefficients with exact
capacity feasibility. For convex maximization through $k$ measurements,
the circulation box projects to a zonotope. Rational generator directions
determine a polynomial number of exposed candidates at fixed $k$.
Their algebraic values are compared approximately and their endpoint
profiles recovered separately. This is the network application of the
classical zonotope approach, with the irrational endpoint and accuracy
accounting made explicit.

\subsection{Independence and convexity do not make every objective easy}

Independent coefficient intervals can realize a circulation cube on a
bounded-data cactus. A coupled convex quadratic on that cube encodes
Max-Cut. Thus convex maximization in all flows is strongly NP-hard,
even with unit terminal nominations, maximum degree three, and resistances
in $[1,16]$; absolute-error $1/4$ approximation remains hard
(\cref{thm:a-design-maxcut}). The positive fixed-measurement theorem
restricts the dimension of the measured image, not merely the degree
or convexity of the objective.

Global correlations can destroy even linear-flow tractability. Two
cycles sharing one resistance already have a curved attainable image:
in \cref{ex:a-corr-shared}, two selected flows satisfy
$x_2=2x_1/(1+x_1)$ on $1/3\le x_1\le1/2$.
Their difference has its strict maximum at the interior shared
resistance $\theta=2$, although its two endpoint values agree. This
example alone proves no hardness, but explains why independent-cycle
endpoint reasoning cannot be reused. A separate paired-triangle
construction proves strong hardness for total-flow maximization and
minimum dissipation over a global resistance polytope, with all
resistances in $[1,3]$, positive flows, and a fixed absolute gap
(\cref{thm:a-corr-hardness}).

For symmetric quadratic laws write
\[
 D(\beta)=b^\top\pi(\beta)=\sum_e\beta_e|x_e(\beta)|^3
          =3\min_{Ax=b}\sum_e\beta_e|x_e|^3/3.
\]
This is constitutive dissipation; the primitive energy is $D/3$ at the
physical state. As a pointwise infimum of linear functions of $\beta$,
$D$ is concave. The direction of optimization is therefore decisive.

\begin{theorem}[Maximum dissipation over a global polytope]
\label{thm:main-energy}
On any connected graph with fixed rational balanced nominations and
symmetric quadratic laws, let resistances range over a nonempty bounded
rational polytope with supplied positive rational lower and upper bounds.
Maximum physical dissipation has an exact formulation with two
second-order cones per edge and a rational admissible profile within
$\eps$ of its maximum computable in polynomial accuracy-bit time.
No operating filters are imposed. On a cactus, rational signed arc
capacities may be added: infeasibility is decided, or the same guarantee
holds with every original capacity satisfied exactly.
\end{theorem}

This is \cref{thm:a-corr-energy-max,cor:a-corr-energy-capacity}.
Cubic energy duality is classical; for the symmetric network model it
is already developed in
\citet{gro2017-algorithmic-results-for-potential-based}.
The contribution here includes the global resistance-polytope formulation,
explicit rational output and precision argument, and original-instance
certificate in \cref{prop:a-corr-certificate}. Concavity of effective
resistance also has classical linear-network precedent
\citep{shannon1956-concavity}. Neither this mechanism nor the theorem
extends by assertion to arbitrary potential filters.

\subsection{Rational witnesses need margins beyond the scalar setting}

Exact capacities beyond cacti can force irrational coefficients even
when every capacity is an ordinary positive-width upper bound. A simple
rank-two series-parallel instance with one uncertain resistance gives
such an example (\cref{prop:a-corr-irrational}). Thus rational feasibility
on a cactus reflects the scalar loop equation, rather than a general
property of rational physical models.

Strict slack restores a quantitative rational witness statement on
any graph. For fixed nonzero nominations put $M=\tfrac12\|b\|_1$.
If two symmetric or directional coefficient profiles have lower bound
$\beta_L>0$, and $q_e$ is the absolute symmetric coefficient change
or the larger of the two directional changes on edge $e$, then
\begin{equation}
 \|x'-x\|_\infty\le\frac{M}{2\beta_L}\sum_e q_e.
 \label{eq:main-sensitivity}
\end{equation}
\Cref{prop:a-corr-sharper} proves this bound through regularized
electrical sensitivity and a limiting argument, including zero flows
and reversals. It is four times smaller than the
uniform estimate in \cref{thm:a-corr-lipschitz}; optimality of its
constant is not claimed. A feasible real profile with supplied rational
capacity slack can consequently be rounded safely. The appendix also
extends witness-size recovery to a rational global coefficient polytope
by intersecting an isolating rational box with that polytope. This is a
conditional existence and rounding guarantee; it is not an algorithm
for locating the initial strictly feasible profile.
